\section{实践视角：一个可维护 Coding Agent 的三层检查表}

\subsection{最容易被忽略的是会话和上下文治理}
Week 2 前半部分用最小 Coding Agent 展示了 tool loop 的核心结构，后半部分再引入 MCP。把两部分放在一起看，会发现一个更实际的工程问题：\textbf{一个能长期稳定工作的 Coding Agent，关键不只在于会不会调用工具，而在于怎么管理上下文、状态和失败恢复。}

\begin{knowledgebox}{三层检查表}
\begin{enumerate}
    \item \textbf{模型层}：系统提示词是否清楚说明任务、边界和工具使用原则；
    \item \textbf{工具层}：工具 schema 是否稳定、返回结果是否适合下一步推理；
    \item \textbf{会话层}：上下文是否会膨胀、是否有 reminder、失败后如何恢复。
\end{enumerate}
\end{knowledgebox}

这也是为什么课程里会强调 context 前置加载、system reminders 和 subagent 分流。这些看似 ``产品技巧'' 的东西，实际上是在弥补当前模型在长期任务上的脆弱性。Agent 不是只要会调用工具就够了，它还必须能在长时间、多步骤、多文件的工作中保持目标一致性。

\subsection{本章小结}
如果把 Week 2 学到的内容真正带回工程实践，一个可靠的 Coding Agent 至少要同时过三道关：模型是否知道该做什么、工具是否适合模型调用、会话是否有足够的上下文治理能力。这三者缺一不可。

\subsection{什么时候不该把问题简单归结为 ``上一个 MCP Server''}
Week 2 还有一个很值得保留的工程判断：不是所有集成问题都应该第一时间被包装成一个新的 MCP Server。协议标准化会降低接入成本，但它不会自动解决权限设计、任务切分和观测性问题。

\begin{warningbox}{MCP 不是系统设计的替代品}
\begin{itemize}
    \item 如果底层 API 本身不稳定，MCP 只会把不稳定更标准地暴露出来；
    \item 如果任务边界不清晰，模型即便能调用工具，也可能在错误的层面上自动化；
    \item 如果缺少日志、重试和回滚机制，协议统一反而可能放大错误传播范围。
\end{itemize}
\end{warningbox}

这也是为什么课程一边讲协议，一边仍然强调 agent loop、context management 与产品约束。真正可靠的 Coding Agent，从来不是 ``工具越多越好''，而是 ``能在明确边界内稳定完成任务''。

